\documentclass[11pt,a4paper]{article}
\usepackage[margin=25mm]{geometry}
\usepackage{amsmath,amssymb,amsthm,booktabs,url}
\newtheorem{lemma}{Lemma}
\title{Central involutions and five exact map counts}
\author{Human Compression project}\date{}
\begin{document}\maketitle
Let $G=O'N$. The fixed primary ATLAS table gives one class of
involutions, with centralizer order 161280. The complete subgroup
roster has one class of subgroups of that order, row 7, and it is
maximal. Hence the members of row 7 are precisely the centralizers
$C_G(t)$ of involutions $t$. We use this borrowed classification
and the previously acquired character evaluations to prove five
further entries of Connor--Leemans, Table~1.

\begin{lemma}\label{unique}
Let $x$ be an element of even order in a finite group $G$.
If the largest power of two dividing $|C_G(x)|$ also divides $|x|$,
then $C_G(x)$ has a unique involution. Consequently any subgroup
containing $x$ lies in at most one subgroup of the form $C_G(t)$,
where $t$ is an involution.
\end{lemma}
\begin{proof}
The Sylow 2-subgroup of the cyclic group $\langle x\rangle$ lies
in the center of $C_G(x)$ and, by hypothesis, is a Sylow 2-subgroup
of $C_G(x)$. It is normal, so Sylow conjugacy makes it the only
Sylow 2-subgroup. Every involution belongs to it, and a cyclic
2-group has only one involution. If $H$ contains $x$ and is contained
in $C_G(t)$, then $t\in C_G(x)$. Thus the involution $t$, and hence
its centralizer, are uniquely determined.
\end{proof}

Here are all the relevant classes in the fixed primary table:
\[
\begin{array}{c|rrrr}
|x|&12&16&20&28\\\hline
\text{classes}&12A&16A,16B,16C,16D&20A,20B&28A,28B\\
|C_G(x)|&36&16&20&28
\end{array}
\]
The hypothesis of Lemma~\ref{unique} holds in each case. The
reader extracts the entire 30-class table from the retained primary
HTML and verifies that every sum of class sizes for a given element
order agrees with the trivial-character vector in the previously
acquired character data. In particular, it distinguishes this table
of $G$ from the following table of $G:2$.

For the five types
\[
          (3,20),\ (5,12),\ (5,20),\ (6,20),\ (10,20),
\]
the character evaluations give zero pairs of the specified orders
with involutory product in twelve of the thirteen maximal classes.
Only row 7 contributes. Every nongenerating pair is contained in a
maximal subgroup, and that subgroup must therefore be an involution
centralizer. One generator has order 12 or 20, so the lemma proves
that this maximal overgroup is unique.

Write $x_L(p,q)$ for the total number of ordered pairs of orders
$p,q$ and product order two in $L$, and $e_G(p,q)$ for the number
which generate $G$. Maximal subgroups are self-normalizing in $G$,
and row 7 has class length 2857239. We have therefore proved
\[
 e_G(p,q)=x_G(p,q)-2857239\,x_{H_7}(p,q).
\]
After division by $|\operatorname{Aut}(G)|=921631011840$, the
exact arithmetic is
\[\begin{array}{c|rrr}
(p,q)&x_G/|\mathrm{Aut}(G)|&2857239x_{H_7}/|\mathrm{Aut}(G)|&e_G/|\mathrm{Aut}(G)|\\\hline
(3,20)&45&1&44\\
(5,12)&212&1&211\\
(5,20)&803&3&800\\
(6,20)&1956&8&1948\\
(10,20)&7227&3&7224\\
\end{array}\]

Automorphisms act freely on generating pairs. The last column
therefore counts the required orbits, with no additional division
for duality. All five values agree with the fixed Table~1.

The centralizer orders, maximal-subgroup classification and character
values are borrowed inputs. The contribution is the general
uniqueness criterion and its application to these counts. The same
criterion also shortens several previously proved local uniqueness
arguments; those certificates remain retained. No new group
acquisition, chirality assertion, principal refutation or
whole-source completion is claimed.

\begin{thebibliography}{9}
\bibitem{CL} T. Connor and D. Leemans, \emph{Algorithmic enumeration
of regular maps}, Ars Math. Contemp. 10 (2016), 211--222.
\bibitem{atlas} R. A. Wilson et al., \emph{ATLAS of Finite Group
Representations}, O'Nan group, conjugacy-class table,
\url{https://brauer.maths.qmul.ac.uk/Atlas/v3/spor/ON/}.
The exact retained HTML edition is bound with the certificate.
\bibitem{roster} T. Connor and D. Leemans, \emph{The subgroup lattice
of $O'N$}, fixed 23-page roster, rows 1 and 7.
\end{thebibliography}
\end{document}
